% Autor szablonu: dr Roman Czapla
% Kompilacja: pdfLaTeX (dwukrotnie, aby zaktualizować spis treści).
% Wymagane logo uczelni: images/logo.png (ścieżka względem pliku głównego).
\documentclass[11pt,a4paper,oneside]{article}
\usepackage[utf8]{inputenc}
\usepackage[T1]{fontenc}
\usepackage[polish,provide=*]{babel}
\usepackage{lmodern,microtype}
\usepackage[margin=25mm,headheight=14pt]{geometry}
\usepackage{xcolor,graphicx,tabularx,booktabs,array,enumitem,fancyhdr,titlesec}
\usepackage[hidelinks]{hyperref}
\definecolor{accent}{HTML}{344B59}
\definecolor{hintcolor}{HTML}{58636A}
\hypersetup{pdfauthor={},pdftitle={Projekt inżynierski - Sprawozdanie końcowe}}
\setlength{\parindent}{0pt}
\setlength{\parskip}{6pt}
\renewcommand{\arraystretch}{1.3}
\setlength{\emergencystretch}{2em}
\setlist[itemize]{leftmargin=*,itemsep=2pt,topsep=3pt}
\titleformat{\section}{\Large\bfseries\color{accent}}{\thesection.}{0.5em}{}
\titleformat{\subsection}{\large\bfseries\color{accent}}{\thesubsection.}{0.5em}{}
\pagestyle{fancy}
\fancyhf{}
\fancyhead[L]{\small Projekt inżynierski}
\fancyhead[R]{\small Sprawozdanie końcowe}
\fancyfoot[C]{\small\thepage}
\renewcommand{\headrulewidth}{0.3pt}
\newcolumntype{Y}{>{\raggedright\arraybackslash}X}

% Podpowiedzi są widoczne podczas wypełniania szablonu.
% Przed oddaniem: uzupełnij pola, zastąp przykłady własną treścią
% i zmień poniższy przełącznik na \podpowiedzifalse.
\newif\ifpodpowiedzi
\podpowiedzitrue
%\podpowiedzifalse
\newcommand{\podpowiedz}[1]{\ifpodpowiedzi\par\begingroup\small\color{hintcolor}\textit{Wskazówka: #1}\par\endgroup\fi}
\newcommand{\pole}[1]{\textnormal{[#1]}}

% Dane dokumentu - uzupełnij w jednym miejscu.
% Ścieżkę logo można zmienić poniżej.
\newcommand{\pliklogo}{images/logo.png}
\newcommand{\logouczelni}{%
  \includegraphics[width=0.8\textwidth,height=40mm,keepaspectratio]{\pliklogo}%
}
\newcommand{\tytulprojektu}{\pole{Tytuł projektu}}
\newcommand{\rokakademicki}{\pole{2026/2027}}
\newcommand{\semestrprojektu}{\pole{zimowy / letni}}
\newcommand{\datadokumentu}{\pole{dd.mm.rrrr}}
\newcommand{\wersjadokumentu}{\pole{1.0}}
% Zespół dwuosobowy tylko za zgodą dyrekcji: zmień na \trzeciaosobafalse.
\newif\iftrzeciaosoba
\trzeciaosobatrue
\newcommand{\studentA}{\pole{Imię i nazwisko osoby 1}}
\newcommand{\studentB}{\pole{Imię i nazwisko osoby 2}}
\newcommand{\studentC}{\pole{Imię i nazwisko osoby 3}}
\newcommand{\albumA}{\pole{numer albumu}}
\newcommand{\albumB}{\pole{numer albumu}}
\newcommand{\albumC}{\pole{numer albumu}}
\newcommand{\prowadzacyA}{\pole{Stopień lub tytuł, imię i nazwisko}}
\newcommand{\prowadzacyB}{\pole{Stopień lub tytuł, imię i nazwisko}}

\begin{document}
\begin{titlepage}
\centering
{\large Uniwersytet Komisji Edukacji Narodowej w Krakowie\par}
{\normalsize Instytut Bezpieczeństwa i Informatyki\par}
\vspace{8mm}
\logouczelni\par
\vspace{10mm}
{\Large\bfseries\color{accent} PROJEKT INŻYNIERSKI\par}
\vspace{3mm}
{\large Sprawozdanie końcowe\par}
\vspace{15mm}
{\LARGE\bfseries\tytulprojektu\par}
\vspace{15mm}
{\small ZESPÓŁ PROJEKTOWY\par}
\vspace{3mm}
\begin{tabular}{ll}
\studentA & Nr albumu: \albumA\\
\studentB & Nr albumu: \albumB
\iftrzeciaosoba\\
\studentC & Nr albumu: \albumC\fi
\end{tabular}
\vspace{12mm}

{\small PROWADZĄCY ZAJĘCIA PROJEKTOWE\par}
\vspace{3mm}
\prowadzacyA\par
\prowadzacyB\par
\vfill
Rok akademicki: \rokakademicki\quad Semestr: \semestrprojektu\par
Data dokumentu: \datadokumentu\quad Wersja: \wersjadokumentu\par
\vspace{6mm}
Kraków\par

\end{titlepage}

\newcommand{\podpisy}{%
\par\vspace{5mm}
\begin{tabularx}{\textwidth}{@{}Y p{0.28\textwidth} p{0.28\textwidth}@{}}
\toprule
Imię i nazwisko & Data & Podpis\\
\midrule
\studentA & \dotfill & \dotfill\\[4mm]
\studentB & \dotfill & \dotfill
\iftrzeciaosoba\\[4mm]
\studentC & \dotfill & \dotfill\fi\\[4mm]
\bottomrule
\end{tabularx}}

\tableofcontents
\podpowiedz{Sprawozdanie podsumowuje cały projekt. Opisz uzyskane rezultaty, zgodność z uzgodnionym zakresem i indywidualny wkład. Szczegóły techniczne możesz wskazać przez odwołania do osobnej dokumentacji. Uzupełnij wszystkie pola, usuń przykłady i wyłącz podpowiedzi przed oddaniem. Dla zespołu dwuosobowego zatwierdzonego przez dyrekcję użyj przełącznika w preambule i usuń opis prac trzeciej osoby.}

\section{Opis projektu}
\podpowiedz{Przedstaw problem, cel, odbiorców oraz finalny rezultat. Wyjaśnij praktyczny i inżynierski charakter projektu. Porównaj rezultat z dokumentem ,,Wybór projektu'', uwzględniając późniejsze uzgodnienia.}
\textbf{Problem i cel:} \pole{krótki opis.}\par
\textbf{Rezultat:} \pole{co powstało i do czego służy.}\par
\textbf{Realizacja założeń:} \pole{spełnione wymagania, uzgodnione zmiany i elementy niezrealizowane.}


\section{Opis zrealizowanych prac}
\podpowiedz{Podsumuj etapy realizacji: koncepcję, dobór metod i narzędzi, wykonanie, integrację oraz weryfikację. Dostosuj opis do charakteru projektu; może dotyczyć aplikacji, animacji, prototypu lub innego rezultatu. Nie kopiuj kolejnych raportów bez opracowania podsumowania.}
\pole{Przebieg prac i najważniejsze decyzje wraz z uzasadnieniem.}

\subsection{Weryfikacja i ograniczenia}
\podpowiedz{Opisz, jak sprawdzono realizację celów i jakość rezultatu. Wskaż testy lub inne sprawdzenia, ich wyniki oraz odwołania do dokumentacji. Podaj znane ograniczenia, problemy i możliwe kierunki rozwoju.}
\pole{Sposób weryfikacji, wyniki i ograniczenia.}

\subsection{Indywidualny wkład członków zespołu}
\podpowiedz{Dla każdej osoby opisz konkretne zadania i rezultaty. Przy pracy wspólnej rozróżnij indywidualne obowiązki. Uwzględnij również koncepcję, integrację, testy, dokumentację i organizację pracy. Opis powinien być zgodny z raportami i oświadczeniem o wkładzie.}
\textbf{\studentA:} \pole{zadania i rezultaty.}\par
\textbf{\studentB:} \pole{zadania i rezultaty.}\par
\iftrzeciaosoba
\textbf{\studentC:} \pole{zadania i rezultaty.}\par
\fi

\subsection{Wykorzystanie AI i materiałów zewnętrznych}
\podpowiedz{Wskaż narzędzia AI i zakres ich zastosowania w całym projekcie. Rozróżnij wspomaganie prac od mechanizmów AI stanowiących część rozwiązania. Opisz weryfikację uzyskanych wyników; dla funkcji AI w projekcie podaj także sposób oceny działania i ograniczenia. Nie wymaga się pełnej historii rozmów ani wszystkich promptów. Jeśli nie korzystano z AI, wpisz tę informację.}
\textbf{AI:} \pole{narzędzia, zastosowania i weryfikacja.}\par
\podpowiedz{Wskaż wykorzystane biblioteki, modele, dane, zasoby i cudze fragmenty kodu, ich pochodzenie i warunki wykorzystania. Odróżnij je od własnego wkładu. Szczegółowe źródła wymień w wykazie na końcu dokumentu.}
\textbf{Materiały zewnętrzne:} \pole{wykorzystane materiały i zakres ich zastosowania.}


\section{Wykaz załączników}
\podpowiedz{Wymień wszystkie oddawane pliki i pakiety wraz z krótkim opisem. Zakres dokumentacji powinien być zgodny z ustaleniami prowadzących. Wskaż kompletny projekt z kodem lub innymi plikami źródłowymi, zasobami i plikami wynikowymi oraz obowiązkowe nagranie. Dokumenty oddaje się w PDF. W razie niewymagania danego rodzaju dokumentacji odnotuj uzgodnienie i uzasadnienie pod tabelą. Dodaj wiersze według potrzeb.}
\begin{tabularx}{\textwidth}{@{}p{0.12\textwidth} Y Y@{}}
\toprule
Oznaczenie & Plik / pakiet & Zawartość i przeznaczenie\\
\midrule
Z1 & \pole{nazwa} & \pole{kompletny projekt: kod lub pliki źródłowe, dane i zasoby.}\\
Z2 & \pole{nazwa} & \pole{dokumentacja projektowa - jeśli wymagana.}\\
Z3 & \pole{nazwa} & \pole{dokumentacja użytkowa - jeśli wymagana.}\\
Z4 & \pole{nazwa} & \pole{dokumentacja instalacyjna - jeśli wymagana.}\\
Z5 & \pole{nazwa nagrania} & \pole{prezentacja działania lub rezultatu oddanej wersji.}\\
\bottomrule
\end{tabularx}
\textbf{Uzgodnienia dotyczące dokumentacji:} \pole{zakres i ewentualne odstąpienie od dokumentów nieadekwatnych do projektu.}

\podpowiedz{Dla projektu programistycznego podaj repozytorium i commit lub tag. Dla innych projektów podaj miejsce przechowywania i oznaczenie wersji plików. Odnośnik nie zastępuje przekazania kompletu plików do archiwizacji.}
\textbf{Repozytorium lub miejsce przechowywania:} \pole{adres.}\par
\textbf{Ostateczna wersja:} \pole{identyfikator wersji i data.}

%\newpage
\section{Oświadczenia}
\subsection{Oświadczenie o wykonaniu projektu i wykorzystanych materiałach}
\podpowiedz{Przed podpisaniem upewnijcie się, że treść jest zgodna z rzeczywistym przebiegiem prac i że ujawniono wykorzystanie AI oraz materiałów zewnętrznych. Oświadczenie podpisują wszyscy członkowie zespołu.}

Oświadczamy, że projekt opisany w niniejszym sprawozdaniu został zrealizowany przez nas pod opieką prowadzących zajęcia projektowe w Instytucie Bezpieczeństwa i Informatyki Uniwersytetu Komisji Edukacji Narodowej w Krakowie.

Przedstawiony opis prac i indywidualnego wkładu odpowiada rzeczywistemu przebiegowi realizacji projektu. Wykorzystanie narzędzi AI oraz materiałów zewnętrznych zostało wskazane w sprawozdaniu i odpowiednich miejscach dokumentacji. Materiałów zewnętrznych nie przedstawiamy jako własnego wkładu.

Potwierdzamy, że zweryfikowaliśmy przekazane rezultaty i rozumiemy rozwiązania należące do zakresu naszych prac. Korzystanie z narzędzi wspomagających nie zwalnia nas z odpowiedzialności za poprawność i jakość oddanego projektu.
\podpisy

%\newpage
\subsection{Oświadczenie o indywidualnym wkładzie}
Oświadczamy, że opis indywidualnego wkładu zawarty w niniejszym sprawozdaniu jest zgodny ze stanem faktycznym i został uzgodniony przez wszystkich członków zespołu.
\podpowiedz{Wskaż konkretne obowiązki poszczególnych osób, a nie wyłącznie procenty. Obszary dostosuj do projektu; przy zadaniach wspólnych określ wkład każdej osoby. Tabela uzupełnia opis prac i nie powinna mu przeczyć.}
\begin{tabularx}{\textwidth}{@{}p{0.26\textwidth} Y Y@{}}
\toprule
Obszar prac & Osoba / osoby & Indywidualny wkład\\
\midrule
Koncepcja i założenia & \pole{osoby} & \pole{obowiązki i rezultaty.}\\
Metody i narzędzia & \pole{osoby} & \pole{obowiązki i rezultaty.}\\
Wykonanie i integracja & \pole{osoby} & \pole{obowiązki i rezultaty.}\\
Weryfikacja / testy & \pole{osoby} & \pole{obowiązki i rezultaty.}\\
Dokumentacja i nagranie & \pole{osoby} & \pole{obowiązki i rezultaty.}\\
Pozostałe prace & \pole{osoby} & \pole{rodzaj prac i wkład.}\\
\bottomrule
\end{tabularx}
\podpisy

\section*{Źródła i materiały referencyjne}
\addcontentsline{toc}{section}{Źródła i materiały referencyjne}
\podpowiedz{Wymień źródła przywołane w sprawozdaniu: publikacje, dokumentację, zbiory danych i inne zasoby. Dla materiałów internetowych podaj tytuł, adres i datę dostępu. Pełny wykaz techniczny może znajdować się w dokumentacji projektowej - wskaż wtedy odpowiednią sekcję. Jeśli nie wykorzystano źródeł zewnętrznych, wpisz tę informację.}
\pole{Wykaz źródeł i materiałów referencyjnych.}
\end{document}
